\chapter{Sampling and the Nyquist Theorem}

Every digital instrument we will ever build -- a microphone feeding a laptop, an
accelerometer feeding a microcontroller, the wrist tracker we are working toward --
faces the same first, unavoidable decision. The world is continuous: voltages,
pressures, accelerations all vary smoothly in time. A computer, by contrast, can
only ever hold a \emph{list of numbers}. Somewhere between the analog world and the
digital one, a continuous wiggle must be turned into a finite stream of samples.

The astonishing fact at the heart of this chapter is that, under one clearly stated
condition, \emph{nothing is lost}. A signal can be chopped into a discrete sequence
of instantaneous readings and then perfectly rebuilt, with no approximation, into
the exact continuous waveform you started with. This is the Sampling Theorem, and the
condition is the Nyquist criterion. When the condition is violated, something
specific and visual goes wrong -- the signal does not merely degrade, it \emph{lies}
to you, pretending to be a different signal entirely. That lie is aliasing, and
learning to see it geometrically is the real payoff of this chapter.

\section{What sampling actually does}
\label{ch18:what-sampling-does}

Start with a continuous-time signal $x(t)$, a real number for every real time $t$.
To sample it, we pick a fixed spacing $T_s$ (the \emph{sampling period}) and record
the value of $x$ at the instants $t = nT_s$ for integer $n$:
\[
  x[n] \;=\; x(nT_s), \qquad n \in \Z .
\]
The reciprocal $f_s = 1/T_s$ is the \emph{sampling rate}, measured in samples per
second (Hz). If we sample at $f_s = 1000$\,Hz we take one reading every
$T_s = 1$\,ms. That is the entire operation: read the signal on a regularly spaced
grid of times and throw everything in between away.

\begin{intuitionbox}
Picture the continuous signal as a smooth pencil line drawn on graph paper, and
sampling as standing only on the vertical grid lines and noting where the pencil
line crosses each one. You end up with a row of dots. The deep question of this
chapter is: given only the dots, can you redraw the pencil line uniquely? Sometimes
yes (Nyquist satisfied), sometimes no (aliasing) -- and the boundary between the two
cases is sharp.
\end{intuitionbox}

\begin{figure}[H]
\centering
\begin{tikzpicture}
\begin{axis}[
  width=12cm, height=5.2cm,
  xlabel={time $t$ (ms)}, ylabel={$x$},
  domain=0:10, samples=160,
  xmin=0, xmax=10, ymin=-1.4, ymax=1.4,
  axis lines=left, ytick={-1,0,1},
  enlargelimits=false,
]
\addplot[sigblue, thick] {sin(deg(2*pi*0.3*x))};
\addplot[specred, only marks, mark=*, mark size=1.6pt, samples=11, domain=0:10]
  {sin(deg(2*pi*0.3*x))};
\end{axis}
\end{tikzpicture}
\caption{A continuous signal (blue) sampled once per millisecond ($f_s=1$\,kHz on this
1\,ms-per-unit axis). The red dots are all the computer ever sees; the blue curve is
the truth we hope to recover.}
\label{ch18:fig-sample-dots}
\end{figure}

\subsection{Sampling as multiplication by an impulse train}
\label{ch18:impulse-train}

To analyze sampling we need a clean mathematical model of "read on a grid." The
cleanest one is multiplication by a \emph{Dirac comb} (also called an impulse train),
\[
  \delta_{T_s}(t) \;=\; \sum_{n=-\infty}^{\infty} \delta(t - nT_s),
\]
a row of infinitely thin, infinitely tall spikes spaced $T_s$ apart, each with unit
area. Multiplying our signal by the comb produces the \emph{sampled signal}
\[
  x_s(t) \;=\; x(t)\,\delta_{T_s}(t)
            \;=\; \sum_{n=-\infty}^{\infty} x(nT_s)\,\delta(t - nT_s),
\]
which is a train of impulses whose areas are the sample values. This idealization
keeps the samples living on the continuous-time axis, so we can take its Fourier
transform and compare it to the spectrum of the original. That comparison is the
whole game, and it produces the spectrum-replication picture that makes everything
else obvious.

\begin{formalbox}
The Dirac comb is not an ordinary function; it is a distribution, justified rigorously
by how it behaves inside integrals: $\int f(t)\,\delta(t-a)\,\dd t = f(a)$. We use it
here only as a bookkeeping device. Real hardware does not produce infinite spikes; it
produces finite numbers held for a finite time. Section~\ref{ch18:reconstruction}
returns to what physical converters actually do.
\end{formalbox}

\section{The spectrum-replication view}
\label{ch18:replication}

Here is where geometry takes over from algebra. The Fourier transform of the Dirac
comb is itself a Dirac comb: a comb with spacing $T_s$ in time transforms into a comb
with spacing $f_s = 1/T_s$ in frequency. And we already know (from the convolution
theorem of the properties chapter) that multiplication in time is convolution in
frequency. Putting these two facts together gives the single most important equation
about sampling:
\[
  X_s(f) \;=\; f_s \sum_{k=-\infty}^{\infty} X\!\left(f - k f_s\right).
\]
In words: \textbf{sampling replicates the spectrum.} Take the original spectrum
$X(f)$, make infinitely many copies of it, shift the copies so they sit centered at
every integer multiple of the sampling rate $f_s$, and add them all up. That sum is
the spectrum of the sampled signal.

\begin{happeningbox}
This is what is \emph{really} happening when you sample. You did not just "lose the
in-between values." In the frequency domain you stamped out a periodic pattern: the
original spectrum, copied endlessly along the frequency axis at spacing $f_s$. A
discrete-time signal always has a periodic spectrum, and the period is exactly the
sampling rate. Every weird sampling phenomenon -- aliasing, the Nyquist limit, the
need for anti-alias filters -- is just a statement about whether those copies stay
out of each other's way or crash into each other.
\end{happeningbox}

\begin{figure}[H]
\centering
\begin{tikzpicture}[scale=1.0]
% --- top: original baseband spectrum ---
\begin{scope}[yshift=2.6cm]
  \draw[->] (-3.6,0) -- (3.8,0) node[right]{$f$};
  \draw[->] (0,-0.1) -- (0,1.3) node[above]{$X(f)$};
  \filldraw[fillsig,draw=sigblue,thick] (-0.8,0) -- (0,1) -- (0.8,0) -- cycle;
  \node[font=\small] at (-0.8,-0.3){$-B$};
  \node[font=\small] at (0.8,-0.3){$B$};
  \node[font=\small,sigblue] at (3.0,1.0){baseband};
\end{scope}
% --- bottom: replicated spectrum, well separated ---
\begin{scope}
  \draw[->] (-3.8,0) -- (4.0,0) node[right]{$f$};
  \draw[->] (0,-0.1) -- (0,1.3) node[above]{$X_s(f)$};
  \foreach \c in {-3,0,3}{
    \filldraw[fillsig,draw=sigblue,thick] (\c-0.6,0) -- (\c,0.75) -- (\c+0.6,0) -- cycle;
  }
  \draw[densely dashed,specred] (1.5,0) -- (1.5,1.1);
  \node[specred,font=\small] at (1.5,1.3){$f_s/2$};
  \node[font=\small] at (3,-0.3){$f_s$};
  \node[font=\small] at (-3,-0.3){$-f_s$};
\end{scope}
\end{tikzpicture}
\caption{Top: the original spectrum is band-limited to $\abs{f}\le B$. Bottom: sampling
copies it at every multiple of $f_s$. Here $f_s > 2B$, so the copies have a clear gap
between them; the red line marks the Nyquist frequency $f_s/2$, which falls inside the
gap. Recovery is possible.}
\label{ch18:fig-replication-good}
\end{figure}

\section{Band-limiting, Nyquist, and the Sampling Theorem}
\label{ch18:sampling-theorem}

The replication picture instantly tells us when sampling is lossless. Suppose the
original signal is \emph{band-limited}: its spectrum is exactly zero above some
highest frequency $B$, so all of its energy lives in $\abs{f} \le B$. Each shifted
copy of the spectrum is then a lump of width $2B$ centered at $kf_s$. Two neighboring
lumps, centered at $0$ and $f_s$, will avoid overlapping precisely when the right edge
of the first ($+B$) sits below the left edge of the second ($f_s - B$):
\[
  B \;<\; f_s - B \qquad\Longleftrightarrow\qquad f_s \;>\; 2B .
\]

\begin{theorem}[Nyquist--Shannon Sampling Theorem]
\label{ch18:thm-sampling}
If a continuous-time signal $x(t)$ contains no frequency components at or above $B$
(i.e.\ $X(f)=0$ for $\abs{f}\ge B$), then $x(t)$ is completely determined by its
samples $x(nT_s)$ taken at a uniform rate $f_s = 1/T_s > 2B$. The original signal can
be reconstructed exactly by
\[
  x(t) \;=\; \sum_{n=-\infty}^{\infty} x(nT_s)\,
             \sinc\!\left(\frac{t - nT_s}{T_s}\right),
  \qquad
  \sinc(u) = \frac{\sin(\pi u)}{\pi u}.
\]
\end{theorem}

Two pieces of vocabulary deserve careful separation, because beginners constantly mix
them up.

\begin{itemize}
\item The \textbf{Nyquist rate} is a property of the \emph{signal}: it is $2B$, the
      minimum sampling rate that will work for a signal whose highest frequency is $B$.
\item The \textbf{Nyquist frequency} (or folding frequency) is a property of the
      \emph{sampler}: it is $f_s/2$, the highest frequency a sampler running at $f_s$
      can represent without ambiguity.
\end{itemize}

The sampling theorem is then a single sentence: sample faster than twice the highest
frequency present, equivalently keep the highest frequency below your Nyquist
frequency, and you lose nothing.

\begin{whybox}
Why exactly twice, and not, say, three times? Because a single sinusoid of frequency
$B$ has spectral lines at $+B$ and $-B$. Sampling copies those lines around every
multiple of $f_s$. The copy of the $-B$ line shifted up by $f_s$ lands at $f_s - B$.
For that intruder to stay above your real $+B$ line you need $f_s - B > B$, i.e.\
$f_s > 2B$. The factor of two is the unavoidable price of frequency having both a
positive and a negative half: every real signal's spectrum is symmetric, so a copy
always comes marching in from the right toward your positive frequencies, and you must
outrun it by a factor of two.
\end{whybox}

\begin{historybox}
The result is named for Harry Nyquist (Bell Labs, 1928), who established the $2B$ rate
in the context of telegraph signaling, and Claude Shannon, who in his 1949 paper
\emph{Communication in the Presence of Noise} stated and proved the reconstruction
form we use today. The Russian engineer Vladimir Kotelnikov published an equivalent
theorem in 1933, and E.\,T.\ Whittaker had studied the interpolation series in 1915,
which is why you sometimes see it called the Whittaker--Kotelnikov--Shannon theorem.
Shannon himself was characteristically modest, calling it "common knowledge in the
communication art."
\end{historybox}

\begin{connectionbox}
The reconstruction formula in Theorem~\ref{ch18:thm-sampling} is an
\emph{orthogonal expansion}, exactly like the Fourier series you met earlier. The
shifted sincs $\sinc((t-nT_s)/T_s)$ form an orthogonal basis for the space of
band-limited signals, and the sample values $x(nT_s)$ are the coordinates of $x$ in
that basis. "A band-limited signal is determined by its samples" is the signal-
processing way of saying "a vector is determined by its coordinates" -- the same inner-
product geometry from the linear-algebra companion, now playing out in function space.
\end{connectionbox}

\section{Aliasing: when the copies collide}
\label{ch18:aliasing}

Now violate the condition. Let the signal contain a frequency $B$ that is \emph{above}
the Nyquist frequency, i.e.\ $B > f_s/2$, equivalently $f_s < 2B$. In the replication
picture the lumps are now too wide for their spacing, and neighboring copies overlap.
Where they overlap they add together, and the sum is hopelessly scrambled: there is no
way, looking only at the overlapped region, to know how much came from the original
lump and how much from its neighbor. That overlap is \emph{aliasing}.

\begin{figure}[H]
\centering
\begin{tikzpicture}
  \draw[->] (-4.2,0) -- (4.4,0) node[right]{$f$};
  \draw[->] (0,-0.1) -- (0,1.4) node[above]{$X_s(f)$};
  % three overlapping triangles
  \foreach \c in {-2.4,0,2.4}{
    \draw[fillsig,draw=sigblue,thick,opacity=0.6,fill opacity=0.35,fill=fillsig]
      (\c-1.5,0) -- (\c,0.9) -- (\c+1.5,0) -- cycle;
  }
  \draw[densely dashed,specred] (1.2,0) -- (1.2,1.15);
  \node[specred,font=\small] at (1.2,1.32){$f_s/2$};
  \node[font=\small] at (2.4,-0.3){$f_s$};
  \node[font=\small] at (-2.4,-0.3){$-f_s$};
  \node[specred,font=\small,align=center] at (3.5,0.7){overlap\\(aliasing)};
  \draw[specred,->] (3.3,0.55) -- (1.9,0.25);
\end{tikzpicture}
\caption{Now $f_s < 2B$: each spectral lump is wider than the spacing $f_s$, so the
copies overlap. Energy that belonged at high frequency reappears \emph{folded} back
below the Nyquist frequency $f_s/2$, masquerading as a lower frequency.}
\label{ch18:fig-replication-bad}
\end{figure}

\subsection{Folding: where the alias actually lands}
\label{ch18:folding}

The most useful single fact about aliasing is the formula for where a too-high
frequency \emph{appears} to be. A pure tone at frequency $f_0$, sampled at $f_s$, is
indistinguishable from a tone at the \emph{apparent} (alias) frequency
\[
  f_{\text{alias}} \;=\; \abs{\,f_0 - f_s\,\operatorname{round}(f_0/f_s)\,},
\]
which always lands in $[0, f_s/2]$. The picture is a sheet of paper folded back and
forth like an accordion at every multiple of $f_s/2$: frequencies above Nyquist do not
vanish, they reflect ("fold") down into the legal band.

\begin{figure}[H]
\centering
\begin{tikzpicture}
\begin{axis}[
  width=12cm, height=5cm,
  xlabel={true frequency $f_0$ (kHz)}, ylabel={apparent $f_{\text{alias}}$ (kHz)},
  domain=0:20, samples=2,
  xmin=0, xmax=20, ymin=0, ymax=6,
  axis lines=left,
  xtick={0,5,10,15,20}, ytick={0,5},
  yticklabels={$0$,$f_s/2$},
  enlargelimits=false,
]
% fs = 10 kHz, fs/2 = 5 kHz. Folded triangle wave: reflects at every multiple of 5.
\addplot[sigblue, thick, domain=0:5,   samples=2] {x};
\addplot[sigblue, thick, domain=5:10,  samples=2] {10-x};
\addplot[sigblue, thick, domain=10:15, samples=2] {x-10};
\addplot[sigblue, thick, domain=15:20, samples=2] {20-x};
\draw[gray, densely dotted] (axis cs:5,0)  -- (axis cs:5,5);
\draw[gray, densely dotted] (axis cs:10,0) -- (axis cs:10,5);
\draw[gray, densely dotted] (axis cs:15,0) -- (axis cs:15,5);
\node[font=\small] at (axis cs:10,5.4) {$f_s$};
\node[font=\small] at (axis cs:15,5.4) {$3f_s/2$};
\end{axis}
\end{tikzpicture}
\caption{The folding map for $f_s = 10$\,kHz. The apparent frequency is a triangle wave
that reflects at every multiple of the Nyquist frequency $f_s/2 = 5$\,kHz. A true tone
at $7$\,kHz reads back as $3$\,kHz, exactly as Example~\ref{ch18:ex-fold} finds.}
\label{ch18:fig-folding}
\end{figure}

\begin{example}[A tone that lies]
\label{ch18:ex-fold}
Sample a $7$\,kHz pure tone at $f_s = 10$\,kHz. The Nyquist frequency is $5$\,kHz, and
$7 > 5$, so we expect trouble. Folding: the nearest multiple of $f_s$ to $7$ is
$10\cdot\operatorname{round}(0.7)=10$, so
$f_{\text{alias}} = \abs{7 - 10} = 3$\,kHz. The samples of the $7$\,kHz tone are
\emph{numerically identical} to the samples of a genuine $3$\,kHz tone -- there is no
test, no algorithm, nothing in the data that can tell them apart. The $7$\,kHz tone has
stolen the identity of a $3$\,kHz tone.
\end{example}

\begin{example}[Two cosines, one set of samples]
\label{ch18:ex-two-cos}
Let $f_s = 8$\,kHz. Consider $\cos(2\pi\cdot 1\,t)$ and $\cos(2\pi\cdot 9\,t)$ (kHz).
At sample times $t = n/8$ ms,
\[
  \cos\!\Big(2\pi\cdot 9 \cdot \tfrac{n}{8}\Big)
  = \cos\!\Big(2\pi\cdot 1 \cdot \tfrac{n}{8} + 2\pi n\Big)
  = \cos\!\Big(2\pi\cdot 1 \cdot \tfrac{n}{8}\Big),
\]
because $2\pi n$ is a whole number of turns. The $9$\,kHz cosine and the $1$\,kHz cosine
produce \emph{exactly} the same samples. This is aliasing reduced to one line of
trigonometry: any two frequencies that differ by a multiple of $f_s$ are
indistinguishable once sampled.
\end{example}

\begin{figure}[H]
\centering
\begin{tikzpicture}
\begin{axis}[
  width=12cm, height=5.4cm,
  xlabel={time $t$ (ms)}, ylabel={amplitude},
  domain=0:1, samples=200,
  xmin=0, xmax=1, ymin=-1.4, ymax=1.6,
  axis lines=left, ytick={-1,0,1},
  legend style={at={(0.5,1.02)},anchor=south,legend columns=3,draw=none,font=\small},
  enlargelimits=false,
]
% high-frequency true signal: 9 cycles over the window
\addplot[specred, thick] {cos(deg(2*pi*9*x))};
\addlegendentry{true 9\,kHz}
% aliased low-frequency: 1 cycle over the window
\addplot[sigblue, thick, densely dashed] {cos(deg(2*pi*1*x))};
\addlegendentry{apparent 1\,kHz}
% sample dots at fs = 8 kHz: t = 0,1/8,...,1
\addplot[black, only marks, mark=*, mark size=1.7pt, samples=9, domain=0:1]
  {cos(deg(2*pi*9*x))};
\addlegendentry{samples}
\end{axis}
\end{tikzpicture}
\caption{The wagon-wheel effect in one picture. Both curves pass through every black
sample dot, so the sampler at $f_s=8$\,kHz cannot tell the true $9$\,kHz tone (solid
red) from the slow $1$\,kHz alias (dashed blue). The fast wiggles between samples are
simply invisible to the sampler.}
\label{ch18:fig-aliasing-time}
\end{figure}

\begin{happeningbox}
The famous \emph{wagon-wheel effect} in old films -- spoked wheels that seem to spin
backward or stand still as the stagecoach speeds up -- is aliasing in pure form. A
movie camera samples the world at about $24$ frames per second. When the wheel's spoke-
passing frequency climbs past $12$\,Hz (the camera's Nyquist frequency), each frame
catches the spokes having rotated \emph{almost} a full spoke-spacing, so they appear to
have crept backward. The wheel is spinning fast; the sampled wheel lies and spins slow.
Your eyes are being shown a perfect alias.
\end{happeningbox}

\begin{trapbox}
"My signal only has energy up to $1$\,kHz, so $f_s = 2$\,kHz is fine." Almost never.
The theorem requires \emph{strictly} greater than $2B$, and it requires the spectrum
to be \emph{exactly} zero at and above $B$. Real signals have tails, and real
anti-alias filters cannot cut infinitely sharply. In practice you sample at $2.2\times$
to $2.5\times$ the highest frequency of interest (audio CD: $44.1$\,kHz for a $20$\,kHz
audio band) to leave the filter room to roll off. Sampling at exactly the Nyquist rate
is a theorem about idealized signals, not engineering advice.
\end{trapbox}

\begin{trapbox}
Sampling exactly \emph{at} a frequency you care about is its own trap. Sample
$\sin(2\pi\cdot 5\,t)$ at $f_s = 10$\,kHz and you can hit it precisely at its zero
crossings: every sample is $0$, and the tone vanishes completely. A frequency sitting
exactly on the Nyquist line is ambiguous in phase and amplitude. This is the deeper
reason the theorem says $f_s > 2B$ and not $f_s \ge 2B$.
\end{trapbox}

\section{Anti-alias filtering: stopping the lie before it happens}
\label{ch18:antialias}

Here is the crucial asymmetry: \textbf{aliasing is irreversible}. Once two frequency
copies have overlapped and added, the sum is a single number per frequency bin, and no
amount of clever processing can unmix them. You cannot remove aliasing in software
after sampling. Therefore the defense must be installed \emph{before} the sampler, in
the analog domain.

That defense is the \emph{anti-alias filter}: an analog low-pass filter placed
immediately ahead of the analog-to-digital converter, designed to remove all energy at
and above the Nyquist frequency $f_s/2$ before sampling ever takes place. If nothing
remains above $f_s/2$, there is nothing to fold down, and the replicated copies cannot
overlap. The filter manufactures the band-limited condition that the theorem demands.

\begin{figure}[H]
\centering
\begin{tikzpicture}[node distance=1.05cm and 1.25cm, every node/.style={font=\small}]
  \node (src)  [draw, rounded corners, fill=blue!4]                    {analog $x(t)$};
  \node (aaf)  [draw, rounded corners, fill=orange!10, right=of src]   {anti-alias LPF};
  \node (adc)  [draw, rounded corners, fill=green!8, right=of aaf]     {ADC @ $f_s$};
  \node (dsp)  [draw, rounded corners, fill=violet!8, right=of adc]    {DSP $x[n]$};
  \draw[-{Stealth}] (src) -- (aaf);
  \draw[-{Stealth}] (aaf) -- (adc);
  \draw[-{Stealth}] (adc) -- (dsp);
  \node[below=0.18cm of aaf, gray, align=center, font=\scriptsize]
    {cut everything\\above $f_s/2$};
\end{tikzpicture}
\caption{The acquisition chain. The anti-alias low-pass filter must sit in the analog
domain, ahead of the converter -- after the ADC it is far too late.}
\label{ch18:fig-chain}
\end{figure}

\begin{figure}[H]
\centering
\begin{tikzpicture}
\begin{axis}[
  width=12cm, height=5cm,
  xlabel={frequency $f$}, ylabel={gain},
  domain=0:3, samples=120,
  xmin=0, xmax=3, ymin=-0.05, ymax=1.25,
  axis lines=left, ytick={0,1}, yticklabels={$0$,$1$},
  xtick={1}, xticklabels={$f_s/2$},
  enlargelimits=false,
]
% ideal brick wall
\addplot[sigblue, thick, dashed, samples=2, domain=0:1] {1};
\addplot[sigblue, thick, dashed, samples=2, domain=1:3] {0};
% realistic rolloff (Butterworth-like)
\addplot[specred, thick] {1/sqrt(1+(x/0.95)^16)};
\draw[densely dotted] (axis cs:1,0) -- (axis cs:1,1);
\node[sigblue, font=\small] at (axis cs:0.45,0.6) {ideal brick wall};
\node[specred, font=\small] at (axis cs:1.95,0.45) {real filter};
\end{axis}
\end{tikzpicture}
\caption{An ideal anti-alias filter is a brick wall (dashed): perfect passband, instant
cutoff at $f_s/2$. Real analog filters (solid) roll off gradually, which is exactly why
we oversample -- to give the skirt room to reach the stopband before the Nyquist line.}
\label{ch18:fig-aaf}
\end{figure}

\begin{whybox}
Why must the anti-alias filter be analog? Because aliasing happens \emph{during}
sampling, the instant the continuous signal is read on the grid. By the time you have
digital samples in memory, the high frequencies have already folded down and are
sitting inside your data wearing the disguise of low frequencies. A digital filter
applied afterward sees only the disguised result and removes nothing genuine. The only
moment to act is while the signal is still continuous and the high frequencies are
still where they belong.
\end{whybox}

\begin{applicationbox}
Oversampling converters (the sigma-delta ADCs in essentially every modern audio chip,
phone, and many IMUs) exploit this cleverly. They sample at a very high rate -- often
$64\times$ or $256\times$ the final rate -- so a gentle, cheap analog filter suffices,
then apply a sharp \emph{digital} low-pass filter and discard most samples
(decimation; see the resampling chapter). Pushing the brick-wall job into the digital
domain, where filters can be made arbitrarily sharp and perfectly linear-phase, is one
of the quiet triumphs of modern DSP. The lesson of this chapter is not violated: the
initial analog filter still runs ahead of the high-rate sampler.
\end{applicationbox}

\section{Reconstruction: rebuilding the continuous signal}
\label{ch18:reconstruction}

Sampling's mirror image is reconstruction: turning the sample sequence back into a
continuous waveform. The replication picture makes the recipe obvious. The spectrum of
the sampled signal is the original baseband lump plus all its copies. To recover the
original, simply \emph{delete every copy and keep the one at baseband} -- that is, apply
an ideal low-pass filter with cutoff $f_s/2$. In the frequency domain reconstruction is
a single multiplication by a rectangle.

What does that rectangle look like back in the time domain? The inverse Fourier
transform of an ideal low-pass rectangle is a sinc. So filtering the impulse train with
an ideal low-pass is the same as replacing each impulse with a sinc and adding them up,
which is precisely the interpolation series of Theorem~\ref{ch18:thm-sampling}:
\[
  x(t) \;=\; \sum_{n} x[n]\,\sinc\!\left(\frac{t - nT_s}{T_s}\right).
\]

\begin{intuitionbox}
The sinc is the perfect interpolating shape. It equals $1$ at its own center and
exactly $0$ at every other sample instant ($t = mT_s$, $m\ne 0$). So in the sum, at each
sample time $t=nT_s$ only that sample's sinc contributes and all the others vanish --
the reconstruction passes exactly through every data point. Between the points, the
overlapping sinc tails conspire to draw in the smoothest possible band-limited curve
that threads all the dots. Linear interpolation "connects the dots" with straight
lines; sinc interpolation connects them with the unique band-limited curve, which is
the genuine original.
\end{intuitionbox}

\begin{figure}[H]
\centering
\begin{tikzpicture}
\begin{axis}[
  width=12cm, height=5cm,
  xlabel={$t/T_s$}, ylabel={$\sinc$},
  domain=-4.2:4.2, samples=200,
  xmin=-4.2, xmax=4.2, ymin=-0.32, ymax=1.15,
  axis lines=left, ytick={0,1},
  xtick={-4,-3,-2,-1,0,1,2,3,4},
  enlargelimits=false,
]
\addplot[sigblue, thick] {sin(deg(pi*x))/(pi*x)};
\addplot[specred, only marks, mark=*, mark size=1.4pt, samples=9, domain=-4:4]
  {sin(deg(pi*x))/(pi*x)};
\draw[gray, densely dotted] (axis cs:-4.2,0) -- (axis cs:4.2,0);
\end{axis}
\end{tikzpicture}
\caption{The reconstruction kernel $\sinc(t/T_s)$. It is $1$ at the origin and zero at
every other integer multiple of $T_s$ (red dots). Centering one such curve on each
sample, scaling by the sample value, and summing gives back the exact continuous signal.}
\label{ch18:fig-sinc}
\end{figure}

\begin{trapbox}
The ideal sinc reconstruction is non-causal and infinitely long -- it needs samples
from the infinite past and future, so no real device implements it exactly. Practical
DACs use a \emph{zero-order hold}: each sample is held flat until the next, producing a
staircase. The staircase has a mild built-in low-pass droop (a $\sinc$ in frequency,
this time from holding rather than from interpolating), which is corrected by a gentle
analog \emph{reconstruction filter} after the DAC. Same physics as the input chain, run
in reverse.
\end{trapbox}

\begin{figure}[H]
\centering
\begin{tikzpicture}
\begin{axis}[
  width=12cm, height=5cm,
  xlabel={time $t/T_s$}, ylabel={amplitude},
  domain=0:8, samples=2,
  xmin=0, xmax=8, ymin=-1.3, ymax=1.5,
  axis lines=left, ytick={-1,0,1},
  xtick={0,1,2,3,4,5,6,7,8},
  legend style={at={(0.5,1.02)},anchor=south,legend columns=2,draw=none,font=\small},
  enlargelimits=false,
]
% true band-limited signal
\addplot[sigblue, thick, domain=0:8, samples=120] {sin(deg(2*pi*x/8))};
\addlegendentry{true signal}
% zero-order-hold staircase: hold each sample value across [n, n+1)
\addplot[quotientorange, thick, const plot mark left, domain=0:8, samples=9]
  {sin(deg(2*pi*x/8))};
\addlegendentry{zero-order hold}
\addplot[specred, only marks, mark=*, mark size=1.6pt, samples=9, domain=0:8]
  {sin(deg(2*pi*x/8))};
\end{axis}
\end{tikzpicture}
\caption{A real DAC reconstructs with a zero-order hold (orange staircase): each sample
value is held flat until the next. The staircase tracks the true band-limited curve
(blue) but introduces a step-shaped error and the mild $\abs{\sinc(fT_s)}$ droop
corrected by the analog reconstruction filter.}
\label{ch18:fig-zoh}
\end{figure}

\begin{example}[Sinc interpolation by hand]
\label{ch18:ex-sinc-interp}
Suppose $T_s = 1$ and the only nonzero samples are $x[0]=2$ and $x[1]=-1$, with all
other samples zero. The reconstructed value halfway between them, at $t=0.5$, is
\[
  x(0.5) = 2\,\sinc(0.5) + (-1)\,\sinc(-0.5)
         = 2\cdot\frac{2}{\pi} - 1\cdot\frac{2}{\pi}
         = \frac{2}{\pi} \approx 0.637 .
\]
Note this is \emph{not} the linear-interpolation midpoint $(2+(-1))/2 = 0.5$. The
band-limited reconstruction lands somewhere different because it is the smoothest curve
consistent with \emph{all} the (here zero) samples, not just the two neighbors.
\end{example}

\section{Choosing a sample rate in practice}
\label{ch18:choosing}

The theorem tells you the floor; engineering tells you where to actually stand. The
procedure is always the same three steps.

\begin{enumerate}
\item Decide the highest frequency $B$ you genuinely need to capture.
\item Choose $f_s$ comfortably above $2B$ -- typically $2.2B$ to $2.5B$ -- to leave the
      anti-alias filter room to roll off.
\item Install an analog low-pass anti-alias filter with cutoff near $B$ and adequate
      attenuation by $f_s/2$.
\end{enumerate}

\begin{table}[H]
\centering
\caption{Typical highest frequencies of interest and the sample rates chosen for them.}
\label{ch18:tab-rates}
\begin{tabular}{@{}llll@{}}
\toprule
Domain & Useful band $B$ & Common $f_s$ & Why \\
\midrule
Telephone speech   & $3.4$\,kHz     & $8$\,kHz         & intelligibility, not fidelity \\
Music (CD)         & $20$\,kHz      & $44.1$\,kHz      & edge of hearing $+$ filter margin \\
Pro audio          & $20$\,kHz      & $48$/$96$\,kHz   & headroom, easier filtering \\
IMU (human motion) & $5$--$25$\,Hz  & $100$--$200$\,Hz & gesture dynamics $+$ margin \\
Surface EMG        & $\sim 500$\,Hz & $1$--$2$\,kHz    & muscle spectral content \\
\bottomrule
\end{tabular}
\end{table}

\begin{goalbox}
For the wrist gesture controller, sample-rate choice splits cleanly by sensor.

\textbf{The IMU (accelerometer + gyroscope).} Voluntary human limb motion -- a flick of
the wrist, a tap, a roll of the forearm -- has almost all of its energy below about
$10$--$20$\,Hz; even sharp impacts rarely carry useful structure past $\sim 50$\,Hz. By
the theorem we would need only $\sim 100$\,Hz, but we sample the IMU at $100$--$200$\,Hz
anyway. The extra rate is not for the motion bandwidth; it is for low latency (a MIDI
note must fire within a few milliseconds of the gesture) and for robust feature
extraction. Sampling an IMU at $1$\,kHz to chase "finger detail" is wasted power and
storage -- there is no finger detail in the wrist's motion to alias in the first place.

\textbf{The EMG / FMG channels.} If we add surface EMG to sense muscle activation,
the story changes completely: useful sEMG energy extends to roughly $400$--$500$\,Hz, so
we must sample at $1$--$2$\,kHz and, critically, place a real analog anti-alias filter
ahead of that ADC. Skipping it lets mains-hum harmonics and high-frequency muscle and
motion noise fold straight down into the band where our gesture features live --
silently corrupting the very signal we are trying to classify.

\textbf{The hard limit.} No amount of sample rate turns a wrist IMU into a finger
tracker. Sampling faster cannot create information the sensor never captured: a device
on the forearm sees the gross motion of the limb, not the individual flexion of each
finger joint. Continuous per-finger pose needs sensors \emph{on the fingers} (rings,
gloves) or muscle-level sensing (EMG/FMG) with a trained classifier -- and even then
yields coarse, discrete finger states, not smooth joint angles. The Nyquist theorem
governs how finely we read each channel; it says nothing about what the channel can see.
\end{goalbox}

\begin{buildbox}
On a typical MCU (Arduino-class or a Pi Pico), the safest pattern is a hardware
timer-driven, fixed-rate sample loop -- never sampling "as fast as the loop happens to
run," which produces jitter that smears your spectrum.

\begin{lstlisting}[language=C]
// Fixed 200 Hz IMU sampling driven by a hardware timer ISR.
// Period = 1/200 s = 5000 us. Keep the ISR tiny: read and stash.
volatile int16_t ax, ay, az, gx, gy, gz;
volatile bool    ready = false;

void onTimer() {                       // fires every 5000 us
  imu_read(&ax,&ay,&az,&gx,&gy,&gz);   // one register burst
  ready = true;                        // hand off to the main loop
}

void loop() {
  if (ready) {
    ready = false;
    process_sample(ax,ay,az,gx,gy,gz); // filtering, features, MIDI
  }
}
\end{lstlisting}

Three field-tested rules:
\begin{itemize}
\item Drive the rate from a timer, not from \texttt{delay()} or loop timing -- jitter is
      sampling on an irregular grid, and it injects broadband noise.
\item Configure the IMU's \emph{own} on-chip low-pass (its DLPF/ODR registers) as your
      anti-alias filter; set its bandwidth below $f_s/2$. Most MEMS IMUs include this --
      use it.
\item Match the ODR (output data rate) to your timer rate. If the chip outputs at
      $1$\,kHz but you read at $200$\,Hz without its internal filter engaged, you alias
      on-chip before the data ever reaches you.
\end{itemize}
\end{buildbox}

\begin{labbox}
A five-minute aliasing demo on any laptop. Generate a tone that sweeps upward in
frequency and listen as it folds back down after crossing Nyquist:
\begin{lstlisting}[language=Python]
import numpy as np, sounddevice as sd
fs = 8000                            # deliberately low rate
t  = np.arange(0, 4, 1/fs)
f  = np.linspace(0, 7000, t.size)    # sweep 0 -> 7 kHz (past 4 kHz Nyquist!)
x  = np.sin(2*np.pi*f*t/2)           # chirp (factor 2: instantaneous freq = f)
sd.play(0.2*x, fs); sd.wait()
\end{lstlisting}
You will hear the pitch rise to $4$\,kHz, then -- impossibly -- come back \emph{down}
again as the true frequency keeps climbing. That descending ghost is the alias of the
ascending real tone, folding at the Nyquist frequency. Now raise \texttt{fs} to
$44100$ and the fold vanishes: the sweep climbs cleanly the whole way.
\end{labbox}

\section{Bandpass sampling (a useful surprise)}
\label{ch18:bandpass}

The theorem is usually stated for baseband signals (energy from $0$ to $B$), giving the
familiar $f_s > 2B$. But the deeper requirement is only that the spectral copies not
overlap -- and if a signal's energy lives in a narrow band \emph{far} from zero, the
copies can be made to interleave neatly even at a rate far below $2 \times$(highest
frequency). This is \emph{bandpass} or \emph{undersampling}, and it is how a radio can
digitize a $100$\,MHz signal of, say, $1$\,MHz bandwidth with an ADC running at only a
few megahertz.

\begin{intuitionbox}
The number that matters is the signal's \emph{bandwidth}, not its highest frequency. A
$1$\,MHz-wide lump centered at $100$\,MHz still only needs its copies spaced so they do
not collide, which can be arranged with $f_s$ a little above $2\,\text{MHz}$, provided
$f_s$ is chosen so the lump lands cleanly in one of the replication slots. Sampling
deliberately aliases the band down to a convenient low frequency, on purpose and
losslessly. It is aliasing turned from foe into tool -- which only makes sense once you
truly see sampling as spectrum replication.
\end{intuitionbox}

\begin{aestheticbox}
Step back and notice the unity. One picture -- copies of the spectrum, stamped out at
spacing $f_s$ -- explains the Nyquist rate (copies must not touch), aliasing (copies
overlapping), the wagon wheel (a copy posing as the original), anti-alias filtering
(narrow the lump so copies fit), reconstruction (keep one copy, delete the rest), and
bandpass sampling (slot a distant lump into a gap). Master that single image and the
entire subject of sampling collapses into one geometric idea you can draw on a napkin.
\end{aestheticbox}

\section*{Summary}
\label{ch18:summary}

\begin{itemize}
\item Sampling reads $x(t)$ on a uniform grid $t=nT_s$; in frequency it
      \emph{replicates} the spectrum at every multiple of $f_s=1/T_s$.
\item \textbf{Sampling theorem:} if $X(f)=0$ for $\abs{f}\ge B$ and $f_s > 2B$, the
      samples determine $x(t)$ exactly. $2B$ is the Nyquist \emph{rate}; $f_s/2$ is the
      Nyquist \emph{frequency}.
\item \textbf{Aliasing} is spectral copies overlapping; a too-high tone $f_0$ folds to
      an apparent frequency in $[0,f_s/2]$ and becomes indistinguishable from it. It is
      irreversible.
\item \textbf{Anti-alias filtering} must be analog and precede the ADC; it band-limits
      the signal so the copies cannot overlap.
\item \textbf{Reconstruction} keeps the baseband copy via an ideal low-pass, equivalent
      to summing shifted, scaled sincs through the data points.
\item In practice, oversample ($2.2$--$2.5\times$). For the wrist tracker:
      $\sim 100$--$200$\,Hz for the IMU, $1$--$2$\,kHz with a real anti-alias filter for
      EMG -- and no sample rate ever conjures finger detail the sensor cannot see.
\end{itemize}

\section{Exercises}
\label{ch18:exercises}

\begin{exercise}
A signal contains frequencies up to $4.5$\,kHz. (a) What is the Nyquist rate? (b) You
sample at $9$\,kHz exactly -- is this safe in practice, and why or why not? (c) Suggest
a practical sample rate and justify it.
\end{exercise}

\begin{exercise}
A $13$\,kHz pure tone is sampled at $10$\,kHz. Using the folding formula, find the
apparent (alias) frequency. Then verify your answer by writing the sampled cosine and
reducing its argument modulo $2\pi$.
\end{exercise}

\begin{exercise}
Two sinusoids of frequency $f_1$ and $f_2$ produce identical samples at rate $f_s$.
Show that $f_1$ and $f_2$ must satisfy $f_2 = f_1 + k f_s$ or $f_2 = -f_1 + k f_s$ for
some integer $k$. Interpret the second family in terms of the folding picture.
\end{exercise}

\begin{exercise}
Sketch, by hand, the replicated spectrum $X_s(f)$ for a triangular baseband spectrum of
half-width $B = 3$\,kHz when (a) $f_s = 10$\,kHz and (b) $f_s = 5$\,kHz. In which case do
the copies overlap, and over what frequency interval?
\end{exercise}

\begin{exercise}
Explain in one or two sentences why an anti-alias filter implemented as a digital FIR
filter applied to the samples cannot prevent aliasing, no matter how sharp it is made.
\end{exercise}

\begin{exercise}
The wagon wheel of a stagecoach has $12$ evenly spaced spokes and is filmed at
$24$\,fps. (a) At what wheel rotation rate (rev/s) does the wheel first appear
\emph{stationary}? (b) Just below that rate, which way does the wheel appear to turn?
\end{exercise}

\begin{exercise}
Using the sinc interpolation series with $T_s=1$, reconstruct the value at $t = 0.25$
for a signal whose only nonzero samples are $x[0]=1$ and $x[2]=1$. Compare with linear
interpolation between the nearest samples.
\end{exercise}

\begin{exercise}
A zero-order-hold DAC outputs each sample held for $T_s$. Its frequency response has
magnitude $\abs{\sinc(f T_s)}$. Compute the gain droop (in dB) at $f = f_s/4$ relative
to DC, and explain why a reconstruction filter is needed.
\end{exercise}

\begin{exercise}[Toward the wrist tracker]
You are logging a wrist IMU for gesture recognition. The fastest deliberate motion you
care about is a wrist flick with energy up to $\sim 15$\,Hz. (a) What is the minimum
sample rate? (b) Why might you still choose $200$\,Hz? (c) Your IMU streams at a fixed
internal rate of $1$\,kHz but you only read every fifth sample in software, with the
chip's on-board low-pass disabled. Describe precisely what goes wrong and how to fix it.
\end{exercise}

\begin{exercise}[Bandpass sampling]
A radio signal occupies $99.5$ to $100.5$\,MHz (bandwidth $1$\,MHz). Explain why an ADC
running at only a few MHz could, in principle, capture it without loss, and state what
property of $f_s$ must be arranged. What real-world impairment makes this harder than
the clean theory suggests?
\end{exercise}
